\subsection{Vector tests behind the slope rules}

Let two lines have nonzero direction vectors $\mathbf v_1$ and $\mathbf v_2$.
Their geometric directions are parallel precisely when one direction vector is a
nonzero scalar multiple of the other:
\[
\boxed{
\mathbf v_1\parallel\mathbf v_2
\iff
\mathbf v_2=\lambda\mathbf v_1
\text{ for some }\lambda\neq0.
}
\]
In three coordinates the same condition will be detected by a zero cross
product. In the plane, the scalar-multiple criterion is already sufficient.

The directions are perpendicular precisely when
\[
\boxed{
\mathbf v_1\perp\mathbf v_2
\iff
\mathbf v_1\cdot\mathbf v_2=0.
}
\]
These tests do not require a slope and therefore include vertical directions
without a separate convention.

For nonvertical lines with slopes $m_1,m_2$, convenient direction vectors are
\[
\mathbf v_1=(1,m_1),
\qquad
\mathbf v_2=(1,m_2).
\]
Their dot product is
\[
\mathbf v_1\cdot\mathbf v_2=1+m_1m_2.
\]
Consequently,
\[
\mathbf v_1\perp\mathbf v_2
\iff
1+m_1m_2=0
\iff
\boxed{m_1m_2=-1}.
\]
Thus the familiar slope rule is not an independent fact. It is the coordinate
form of the dot-product test for the special direction vectors $(1,m)$.

\subsection{The nearest point on a line}

Let
\[
L=\{P_0+t\mathbf v:t\in\mathbb R\},
\qquad \mathbf v\neq\mathbf0,
\]
and let $P$ be a point not necessarily on $L$. The displacement
\[
P-P_0
\]
can be decomposed into a component parallel to $\mathbf v$ and a component
perpendicular to $\mathbf v$. The parallel component is
\[
\operatorname{proj}_{\mathbf v}(P-P_0)
=
\frac{(P-P_0)\cdot\mathbf v}{\mathbf v\cdot\mathbf v}\,\mathbf v.
\]
Starting from $P_0$ and moving by this component reaches the foot of the
perpendicular:
\[
\boxed{
Q=P_0+\operatorname{proj}_{\mathbf v}(P-P_0).
}
\]
The point $Q$ lies on the line, and $P-Q$ is perpendicular to the line. It is
therefore the shortest displacement from $P$ to any point of $L$.

\begin{center}
\begin{tikzpicture}[scale=0.95,>=stealth]
    \coordinate (P0) at (0.8,0.8);
    \coordinate (Q) at (3.7,2.1);
    \coordinate (P) at (2.8,4.1);
    \coordinate (R) at (5.4,2.86);
    \draw[subset] (0.1,0.49) -- (R) node[above right] {$L$};
    \draw[blue!70!black,line width=1.1pt,->] (P0) -- (Q)
        node[midway,below right] {$\operatorname{proj}_{\mathbf v}(P-P_0)$};
    \draw[orange!85!black,line width=1.1pt,->] (Q) -- (P)
        node[midway,right] {$P-Q$};
    \draw[very thick,->] (P0) -- (P)
        node[midway,above left] {$P-P_0$};
    \draw (Q) ++(-0.22,-0.10) -- ++(-0.12,0.27) -- ++(0.22,0.10);
    \fill[geometricSubsetColor] (P0) circle (2pt) node[below left] {$P_0$};
    \fill[geometricSubsetColor] (Q) circle (2pt) node[below right] {$Q$};
    \fill[geometricSubsetColor] (P) circle (2pt) node[above] {$P$};
\end{tikzpicture}
\end{center}

Hence
\[
\boxed{
d(P,L)
=
\left\|
(P-P_0)-\operatorname{proj}_{\mathbf v}(P-P_0)
\right\|.
}
\]

\subsection{Distance from the normal equation}

Suppose the same line is written as
\[
ax+by=c,
\qquad (a,b)\neq(0,0),
\]
with normal vector
\[
\mathbf n=(a,b).
\]
For $P=(x_P,y_P)$ and any fixed $P_0$ on the line,
\[
\mathbf n\cdot(P-P_0)=ax_P+by_P-c.
\]
The signed scalar component of $P-P_0$ in the normal direction is therefore
\[
\frac{ax_P+by_P-c}{\sqrt{a^2+b^2}}.
\]
Distance is nonnegative, so
\[
\boxed{
d(P,L)
=
\frac{|ax_P+by_P-c|}{\sqrt{a^2+b^2}}.
}
\]
This formula is not a new rule detached from projection. It is the length of the
normal component of the displacement from the line to the point.
